\section{一阶形变}\label{src:C.1}
\begin{epigraph}
园丁敏锐地留意自然秩序的深层模式, 同时也知道, 总会有许多事物超出他们的视野.
\epigraphauthor{James Carse, \textit{有限与无限的游戏}}
\end{epigraph}
对域 $\Bbbk$, 将\newterm{对偶数}{dual numbers} 记为 $\Bbbk[\epsilon]:=\Bbbk[\epsilon]/(\epsilon^2)$.
\subsection{一阶嵌入形变}\label{src:C.1.1-subsection}
\begin{definition}\label{src:C.1.1}
设 $X$ 为 $\Bbbk$ 上的射影概形, $Z_0\subseteq X$ 为闭子概形. $Z_0\subseteq X$ 的\newterm{一阶形变}{first-order deformation} 是闭子概形 $Z\subseteq X_{\Bbbk[\epsilon]}:=X\times_{\Bbbk}\Bbbk[\epsilon]$, 它在 $\Bbbk[\epsilon]$ 上平坦且满足 $Z_0=Z\times_{\Bbbk[\epsilon]}\Bbbk$. 图示而言, 一阶形变是在图
\[
\begin{tikzcd}[column sep=small]
&X\arrow[rr,hook]\arrow[dd]&&X_{\Bbbk[\epsilon]}\arrow[dd]\\
Z_0\arrow[ur,hook,"\text{闭}"]\arrow[dr]\arrow[rr,dashed]&&Z\arrow[ur,dashed,hook,"\text{闭}"']\arrow[dr,dashed,"\text{平坦}"]\\
&\Spec\Bbbk\arrow[rr,hook]&&\Spec\Bbbk[\epsilon]
\end{tikzcd}
\]
中补入在 $\Bbbk[\epsilon]$ 上平坦的概形 $Z$ 以及虚线箭头, 使图为 Cartesian 图.

注意 $Z_0$ 和 $Z$ 的中心纤维 $Z\times_{\Bbbk[\epsilon]}\Bbbk$ 都嵌入 $X$, 因此定义中要求它们相等是有意义的. 若 $Z=Z_0\times_{\Bbbk}\Bbbk[\epsilon]$, 则称 $Z\subseteq X_{\Bbbk[\epsilon]}$ 为\textit{平凡的}.
\end{definition}
\begin{remark}\label{src:C.1.2}
闭子概形 $Z_0\subseteq X$ 定义 Hilbert 概形的 $\Bbbk$-点 $[Z_0\subset X]\in\Hilb^P(X)$, 其中 $P$ 为 $Z_0$ 关于 $X$ 上某条固定丰沛线丛的 Hilbert 多项式. 一阶形变对应交换图
\[
\begin{tikzcd}
\Spec\Bbbk\arrow[r,"{[Z_0\subseteq X]}"]\arrow[d,hook]&\Hilb^P(X)\\
\Spec\Bbbk[\epsilon]\arrow[ur,"{[Z\subseteq X_{\Bbbk[\epsilon]}]}"']&
\end{tikzcd}
\]
亦即切向量 $[Z\subseteq X_{\Bbbk[\epsilon]}]\in T_{\Hilb^P(X),[Z_0\subseteq X]}$.
\end{remark}
\begin{proposition}\label{src:C.1.3}
设 $X$ 为 $\Bbbk$ 上的概形, $Z_0\subseteq X$ 为由理想层 $I_0\subseteq\mathcal O_X$ 定义的闭子概形. 有双射
\[
\{\text{一阶形变 }Z\subseteq X_{\Bbbk[\epsilon]}\}\simeq\mathrm H^0(Z_0,N_{Z_0/X}),
\]
其中 $N_{Z_0/X}=\mathcal Hom_{\mathcal O_{Z_0}}(I_0/I_0^2,\mathcal O_{Z_0})$ 为法层. 此对应将平凡形变对应到 $0\in\mathrm H^0(Z_0,N_{Z_0/X})$.
\end{proposition}
\begin{remark}\label{src:C.1.4}
由注记 \bookref{C.1.2}, 此命题给出双射 $T_{\Hilb^P(X),[Z_0\subseteq X]}\simeq\mathrm H^0(Z_0,N_{Z_0/X})$.
\end{remark}
\begin{proof}
先处理 $X=\Spec B$ 且 $Z_0=\Spec B/I_0$ 的情形, 证明一阶形变集与
\[
\mathrm H^0(Z_0,N_{Z_0/X})\simeq\operatorname{Hom}_{B/I_0}(I_0/I_0^2,B/I_0)\simeq\operatorname{Hom}_B(I_0,B/I_0)
\]
双射. 给定一阶形变 $Z=\Spec B[\epsilon]/I$, $Z$ 在 $\Bbbk[\epsilon]$ 上的平坦性意味着, 将 $\Bbbk[\epsilon]$-模正合列 $0\to I\to B[\epsilon]\to B[\epsilon]/I\to0$ 与 $\Bbbk=\Bbbk[\epsilon]/(\epsilon)$ 作张量积, 得到正合列 $0\to I_0\to B\to B/I_0\to0$. 如下定义 $\alpha:I_0\to B/I_0$: 对 $x_0\in I_0$, 选取原像 $x=a+b\epsilon\in I$, 并令 $\alpha(x_0):=\overline b\in B/I_0$. 反之, 给定 $B$-模同态 $\alpha:I_0\to B/I_0$, 定义
\[
I=\{a+b\epsilon\mid a\in I_0,\ b\in B,\ \overline b=\alpha(a)\in B/I_0\}\subseteq B[\epsilon].
\]
则 $(B[\epsilon]/I)\otimes_{\Bbbk[\epsilon]}\Bbbk=B/I_0$. 为证明 $B[\epsilon]/I$ 在 $\Bbbk[\epsilon]$ 上平坦, 由双数环上的平坦性判别法 (见注记 \bookref{A.2.7}), 只需检验映射 $B/I_0\xrightarrow{\epsilon}B[\epsilon]/I$ 单射: 给定 $b\in B$ 满足 $\epsilon b\in I$, 则由 $I$ 的定义有 $b\in I_0$. 因此 $Z=\Spec B[\epsilon]/I$ 定义 $Z_0$ 的一阶形变. 对一般 $\Bbbk$-概形 $X$, 选取仿射覆盖 $\{U_i\}$ 后, 可检验 $U_i$ 的形变与 $\mathrm H^0(U_i\cap Z_0,N_{U_i\cap Z_0/U_i})$ 之间的双射粘合为所需双射. 详情另见 \cite[定理 $6.1$]{Art76}, \cite[定理 $2.8$]{Kol96}, \cite[命题 $3.2.1$]{Ser06}, \cite[命题 $2.3$]{Har10}.
\end{proof}
\subsection{概形的一阶形变}\label{src:C.1.2-subsection}
\begin{definition}\label{src:C.1.5}
设 $X_0$ 为域 $\Bbbk$ 上的概形. $X_0$ 的一阶形变是在 $\Bbbk[\epsilon]$ 上平坦的概形 $X$, 连同同构 $\alpha:X_0\to X\times_{\Bbbk[\epsilon]}\Bbbk$, 亦即 Cartesian 图
\begin{equation}\label{src:C.1.6}
\begin{tikzcd}
X_0\arrow[r,hook,dashed]\arrow[d]&X\arrow[d,dashed,"\text{平坦}"]\\
\Spec\Bbbk\arrow[r,hook]\arrow[ur,phantom,"\square",very near start]&\Spec\Bbbk[\epsilon].
\end{tikzcd}
\end{equation}
一阶形变 $(X,\alpha)$ 与 $(X',\alpha')$ 之间的\textit{态射} 是 $\Bbbk[\epsilon]$ 上的概形态射 $\beta:X\to X'$, 满足 $(\beta\times_{\Bbbk[\epsilon]}\Bbbk)\circ\alpha=\alpha'$; 换言之, 将 $X$ 和 $X'$ 放入 Cartesian 图 \bookeqref{C.1.6} 时, 要求 $\beta$ 在中心纤维 $X_0$ 上的限制为恒等态射.

若 $X$ 与 $X_0\times_{\Bbbk}\Bbbk[\epsilon]$ 作为一阶形变同构, 则称 $X$ 为\textit{平凡的}; 若存在 Zariski 覆盖 $X=\bigcup_iU_i$, 使得 $U_i$ 为 $U_i\times_{\Bbbk[\epsilon]}\Bbbk\subseteq X_0$ 的平凡一阶形变, 则称其为\textit{局部平凡的}.
\end{definition}
每个形变态射必为同构. 这是下述代数事实的推论.
\begin{lemma}\label{src:C.1.7}
设 $A$ 为环, $\mathfrak m\subseteq A$ 为幂零理想 (例如 $(A,\mathfrak m)$ 为 Artin 局部环), $M\to N$ 为 $A$-模同态. 假设 $N$ 在 $A$ 上平坦. 若 $M/\mathfrak mM\to N/\mathfrak mN$ 为同构, 则 $M\to N$ 也是同构.
\end{lemma}
\begin{proof}
若 $C:=\operatorname{coker}(M\to N)$, 则 $C/\mathfrak mC=\operatorname{coker}(M/\mathfrak mM\to N/\mathfrak mN)=0$. 由于存在 $n$ 使 $\mathfrak m^n=0$, 得 $C=\mathfrak mC=\mathfrak m^2C=\cdots=\mathfrak m^nC=0$. 若 $K:=\ker(M\to N)$, 则 $N$ 的平坦性给出 $K/\mathfrak mK=\ker(M/\mathfrak mM\to N/\mathfrak mN)=0$. 因此 $K=\mathfrak mK=\cdots=\mathfrak m^nK=0$, 从而 $M\to N$ 为同构.
\end{proof}
\begin{proposition}\label{src:C.1.8}
域 $\Bbbk$ 上光滑仿射概形 $X_0$ 的每个一阶形变都平凡. 换言之, $X_0$ 是\newterm{刚性的}{rigid}.
\end{proposition}
\begin{proof}
设 $X$ 为 $X_0$ 的一阶形变. 因为 $X_0\to\Spec\Bbbk$ 光滑, 可应用光滑性的无穷小提升判据 (\bookref{A.3.1}) 构造提升 $X\to X_0$, 使图
\[
\begin{tikzcd}
X_0\arrow[r,"\mathrm{id}"]\arrow[d,hook]&X_0\arrow[d,"\text{光滑}"]\\
X\arrow[r]\arrow[ur,dashed]&\Spec\Bbbk
\end{tikzcd}
\]
交换. 它诱导 $\Bbbk[\epsilon]$ 上的态射 $X\to X_0\times_{\Bbbk}\Bbbk[\epsilon]$, 其在 $X_0$ 上的限制为恒等态射, 因此由引理 \bookref{C.1.7} 为同构. 另见 \cite[习题 II.$8.7$]{Har77}.
\end{proof}
若 $X_0$ 不光滑或不仿射, 则一阶形变未必平凡. 例如, 若 $X_0=\Spec\Bbbk[x,y]/(xy)$ 为结点仿射平面曲线, 则 $X=\Spec\Bbbk[x,y,\epsilon]/(xy-\epsilon)\to\Spec\Bbbk[\epsilon]$ 为非平凡一阶形变. 另一方面, 考虑 $\lambda\neq0,1$ 时的椭圆曲线 $E_0=V(y^2z-x(x-z)(x-\lambda z))\subseteq\mathbb P^2$,
\begin{equation}\label{src:C.1.9}
E_\alpha=V(y^2z-x(x-z)(x-(\lambda+\alpha\epsilon)z))\subseteq\mathbb P^2_{\Bbbk[\epsilon]}
\end{equation}
对每个 $\alpha\in\Bbbk$ 定义 $E_0$ 的一阶形变, 而 $\alpha\mapsto E_\alpha$ 定义 $\Bbbk$ 与一阶形变同构类集合之间的双射; 这由命题 \bookref{C.1.11} 推出, 因为 $\mathrm H^1(E_0,T_{E_0})=\mathrm H^1(E_0,\mathcal O_{E_0})=\Bbbk$. 实际上, 同一公式 \bookeqref{C.1.9} 定义 $E_0$ 在 $\mathbb A^1$ 上的整体形变, 带有两条奇异纤维.

为理解局部平凡形变, 我们需要了解平凡形变的自同构.
\begin{lemma}\label{src:C.1.10}
设 $X_0=\Spec A$ 为 $\Bbbk$ 上的仿射概形, $X=\Spec A[\epsilon]$ 为平凡一阶形变. 有等同
\[
\{\text{一阶形变的自同构 }X\to X\}\simeq\operatorname{Der}_{\Bbbk}(A,A)\simeq\operatorname{Hom}_A(\Omega_{A/\Bbbk},A).
\]
\end{lemma}
\begin{proof}
第二个等同由微分模的万有性质给出. $X$ 的自同构 (作为一阶形变) 对应于模 $\epsilon$ 为恒等映射的 $\Bbbk[\epsilon]$-代数同构 $\phi:A\oplus A\epsilon\to A\oplus A\epsilon$. 因此 $\phi$ 由元素 $a\in A$ 的像 $\phi(a)=a+d(a)\epsilon$ 决定, 其中 $d:A\to A$ 为 $\Bbbk$-线性映射. 映射 $\phi$ 为环同态, 当且仅当对元素 $a,a'\in A$ 有
\[
aa'+d(aa')\epsilon=(a+d(a)\epsilon)(a'+d(a')\epsilon)=aa'+(ad(a')+a'd(a))\epsilon;
\]
这恰好是 $d:A\to A$ 为 $\Bbbk$-导子的条件.
\end{proof}
对 $\Bbbk$ 上的概形 $X_0$, 用 $\operatorname{Def}(X_0)$ 和 $\operatorname{Def}^{\mathrm{lt}}(X_0)$ 分别表示一阶形变及局部平凡一阶形变的同构类集合.
\begin{proposition}\label{src:C.1.11}
对具有仿射对角线的有限型 $\Bbbk$-概形 $X_0$, 有双射
\[
\operatorname{Def}^{\mathrm{lt}}(X_0)\simeq\mathrm H^1(X_0,T_{X_0}),
\]
其中 $T_{X_0}=\mathcal Hom_{\mathcal O_{X_0}}(\Omega_{X_0/\Bbbk},\mathcal O_{X_0})$. 平凡形变对应 $0\in\mathrm H^1(X_0,T_{X_0})$. 若 $X_0$ 还在 $\Bbbk$ 上光滑, 则有双射
\[
\operatorname{Def}(X_0)\simeq\mathrm H^1(X_0,T_{X_0}).
\]
\end{proposition}
\begin{proof}
设 $X\to\Spec\Bbbk[\epsilon]$ 为 $X_0$ 的局部平凡一阶形变. 选取 $X_0$ 的仿射覆盖 $\{U_i\}$ 及同构 $\phi_i:U_i\times_{\Bbbk}\Bbbk[\epsilon]\xrightarrow{\sim}X\cap U_i$. 令 $U_{ij}=U_i\cap U_j$, 则有平凡形变 $U_{ij}\times_{\Bbbk}\Bbbk[\epsilon]$ 的自同构 $\phi_j^{-1}|_{U_{ij}\times_{\Bbbk}\Bbbk[\epsilon]}\circ\phi_i|_{U_{ij}\times_{\Bbbk}\Bbbk[\epsilon]}$, 由引理 \bookref{C.1.10}, 它对应元素 $\phi_{ij}\in T_{X_0}(U_{ij})$. 因为在 $U_{ijk}:=U_i\cap U_j\cap U_k$ 上有 $\phi_{ij}\circ\phi_{jk}=\phi_{ik}$, 得 $\phi_{ij}+\phi_{jk}=\phi_{ik}\in T_{X_0}(U_{ijk})$. 回顾 $\mathrm H^1(X_0,T_{X_0})$ 可用 Čech 复形计算:
\[
\begin{aligned}
\bigoplus_iT_{X_0}(U_i)&\xrightarrow{d_0}\bigoplus_{i,j}T_{X_0}(U_{ij})\xrightarrow{d_1}\bigoplus_{i,j,k}T_{X_0}(U_{ijk}),\\
(s_{ij})&\longmapsto(s_{ij}|_{U_{ijk}}-s_{ik}|_{U_{ijk}}+s_{jk}|_{U_{ijk}})_{ijk}.
\end{aligned}
\]
因为 $\{\phi_{ij}\}\in\bigoplus_{i,j}T_{X_0}(U_{ij})$ 属于 $\ker d_1$, 它定义 $\mathrm H^1(X_0,T_{X_0})=\ker d_1/\im d_0$ 的元素. 反之, 给定 $\mathrm H^1(X_0,T_{X_0})$ 的元素并选取代表 $\{\phi_{ij}\}\in\ker d_1$, 将每个 $\phi_{ij}$ 看成 $U_{ij}$ 的平凡形变的自同构, 可沿 $U_{ij}\times_{\Bbbk}\Bbbk[\epsilon]$, 经由 $\phi_{ij}$ 粘合平凡形变 $U_i\times_{\Bbbk}\Bbbk[\epsilon]$, 构造 $X_0$ 的整体一阶形变 $X$. 最后的断言由 $X_0$ 光滑时所有一阶形变都局部平凡 (命题 \bookref{C.1.8}) 推出. 另见 \cite[习题 III.$4.10$ 和例 III.$9.13.2$]{Har77}.
\end{proof}
\begin{example}\label{src:C.1.12}
若 $C$ 为亏格 $g\geq2$ 的光滑射影曲线, 则
\[
T_{\mathcal M_g,[C]}=\mathrm H^1(C,T_C)\overset{\mathrm{SD}}{=}\mathrm H^0(C,\Omega_{C/\Bbbk}^{\otimes2}),
\]
由 Riemann--Roch, 这是 $3g-3$ 维向量空间.
\end{example}
\begin{exercise}[简单]\label{src:C.1.13}
用 Euler 正合列证明 $\mathrm H^1(\mathbb P^n,T_{\mathbb P^n})=0$, 并推出 $\mathbb P^n$ 的每个一阶形变都平凡, 即 $\mathbb P^n$ 是刚性的.
\end{exercise}
\begin{exercise}[中等]\label{src:C.1.14}
设 $C_0$ 为 $\Bbbk$ 上的光滑固有曲线, 带有标记点 $p_1,\ldots,p_n\in C_0(\Bbbk)$. 用 $\operatorname{Def}(C_0,p_i)$ 表示 $(C_0,p_i)$ 的一阶形变集, 即带有延拓 $(C_0,p_i)$ 的 $n$ 个截面 $\sigma_i$ 的平坦态射 $C\to\Spec\Bbbk[\epsilon]$. 证明
\[
\operatorname{Def}(C_0,p_i)\simeq\mathrm H^1\left(C_0,T_{C_0}(-\sum_i p_i)\right).
\]
\end{exercise}
\begin{remark}\label{src:C.1.15}
更一般地, 若 $X_0$ 在 $\Bbbk$ 上泛光滑且为局部完全交, 则
\[
\operatorname{Def}(X_0)=\operatorname{Ext}_{\mathcal O_{X_0}}^1(\Omega_{X_0},\mathcal O_{X_0}).
\]
注意 $X_0$ 光滑时, 这重新给出命题 \bookref{C.1.11}. 若 $X_0=\Spec R$ 为带有孤立奇点的完备 Noether 局部环的谱, 且 $R$ 为完全交环, 这也成立. 容易看出 $X_0$ 的一阶形变 $X$ 如何诱导扩张类: 闭浸入 $X_0\hookrightarrow X$ 由理想 $J\subset\mathcal O_X$ 定义, 满足 $J^2=0$ 且作为 $\mathcal O_{X_0}$-模有 $J\simeq\mathcal O_{X_0}$, 而通常的右正合列
\[
\mathcal O_{X_0}\to\Omega_X|_{X_0}\to\Omega_{X_0}\to0
\]
也左正合, 因为最左映射的核一方面无挠 (作为 $\mathcal O_{X_0}$ 的子层), 另一方面又是挠的 (因为泛光滑性意味着该映射泛单射).
\end{remark}
\begin{exercise}[适合练习]\label{src:C.1.16}
\begin{enumerate}[(1)]
\item 证明有双射 $\operatorname{Def}(\Bbbk[x,y]/(xy))\simeq\Bbbk$, 其中元素 $t\in\Bbbk$ 对应一阶形变 $\Spec\Bbbk[x,y,\epsilon]/(xy-t\epsilon)$.
\item 分类 $A_k$-奇点 $\Bbbk[x,y]/(y^2-x^{k+1})$ 的一阶形变.
\end{enumerate}
\end{exercise}
\subsection{向量丛与凝聚层的一阶形变}\label{src:C.1.3-subsection}
\begin{definition}\label{src:C.1.17}
设 $X$ 为 $\Bbbk$ 上的概形, $E_0$ 为凝聚层. $E_0$ 的一阶形变是二元组 $(E,\alpha)$, 其中 $E$ 为 $X\times_{\Bbbk}\Bbbk[\epsilon]$ 上在 $\Bbbk[\epsilon]$ 上平坦的凝聚层, $\alpha:E_0\xrightarrow{\sim}E|_X$ 为同构. 图示为
\[
\begin{tikzcd}
E_0\arrow[d,no head]&E\arrow[d,no head,"{\text{平坦}/\Bbbk[\epsilon]}"]\\
X\arrow[r,hook]&X_{\Bbbk[\epsilon]}.
\end{tikzcd}
\]
一阶形变的态射 $(E,\alpha)\to(E',\alpha')$ 是 $X'$ 上凝聚层的态射 $\beta:E\to E'$ (由引理 \bookref{C.1.7}, 等价地为同构), 满足 $\alpha'=\beta|_X\circ\alpha$. 若 $(E,\alpha)$ 作为一阶形变与 $(p^*E_0,\mathrm{id})$ 同构, 其中 $p:X_{\Bbbk[\epsilon]}\to X$, 则称其为\textit{平凡的}.
\end{definition}
\begin{proposition}\label{src:C.1.18}
设 $X$ 为 $\Bbbk$ 上的概形, $E_0$ 为凝聚层. 有双射
\[
\{E_0\text{ 的一阶形变 }(E,\alpha)\}/{\sim}\ \simeq\operatorname{Ext}_{\mathcal O_X}^1(E_0,E_0).
\]
此对应将平凡形变对应到 $0\in\operatorname{Ext}_{\mathcal O_X}^1(E_0,E_0)$. 若 $E_0$ 还是向量丛 (相应地, 线丛), 则每个一阶形变都是向量丛 (相应地, 线丛), 且 $E_0$ 的一阶形变同构类集合与 $\mathrm H^1(X,\mathcal End_{\mathcal O_X}(E_0))$ (相应地, $\mathrm H^1(X,\mathcal O_X)$) 双射.
\end{proposition}
\begin{proof}
若 $(E,\alpha)$ 为一阶形变, 则由 $E$ 在 $\Bbbk[\epsilon]$ 上平坦, 可将 $\Bbbk[\epsilon]$-模正合列 $0\to\Bbbk\xrightarrow{\epsilon}\Bbbk[\epsilon]\to\Bbbk\to0$ 与 $E$ 作张量积, 得到正合列 $0\to E_0\xrightarrow{\epsilon}E\to E_0\to0$ (经由 $\alpha$ 将 $E\otimes_{\Bbbk[\epsilon]}\Bbbk$ 与 $E_0$ 等同). 因为 $\operatorname{Ext}_{\mathcal O_X}^1(E_0,E_0)$ 参数化扩张的同构类 \cite[习题 III.$6.1$]{Har77}, 这就构造了 $\operatorname{Ext}_{\mathcal O_X}^1(E_0,E_0)$ 的元素. 反之, 给定正合列 $0\to E_0\xrightarrow{\alpha}E\to E_0\to0$, 则 $E$ 为 $X_{\Bbbk[\epsilon]}$ 上的凝聚层, 且由双数环上的平坦性判别法 (见注记 \bookref{A.2.7}), 它在 $\Bbbk[\epsilon]$ 上平坦. 限制 $E|_X$ 经由 $\alpha$ 与 $E_0$ 同构. 另见 \cite[定理 $2.7$]{Har10}.
\end{proof}
\begin{remark}\label{src:C.1.19}
命题 \bookref{C.1.3}, \bookref{C.1.11} 和 \bookref{C.1.18} 的分类赋予一阶形变同构类集合以向量空间结构. 这一结构也可看作 Rim--Schlessinger 齐性条件的推论; 见引理 \bookref{C.4.9}. 在每种情形中, 数乘与加法都有几何描述.
\end{remark}
